% =====================================================================
\chapter{Модели: создание и основные настройки}
\label{ch:models}
% =====================================================================
\chapterintro
  {что такое «модель» в пульте, как её создать и выбрать, как настроить таймеры, триммеры
   и предполётные проверки}
  {пульт после первоначальной настройки}

\section{Что такое модель в пульте}

\begin{simple}
«Модель» в пульте --- это \textbf{листок с настройками} для одного конкретного самолёта или
коптера. В пульте можно хранить десятки таких листков и переключаться между ними, как между
сохранениями в игре. Правило: \textbf{одна настоящая модель --- один листок в пульте}.
\end{simple}

\begin{figure}[H]
\centering
\begin{tikzpicture}[every node/.style={font=\sffamily\small}]
  \pic[transform shape,scale=1.2] at (0,0) {minitx};
  \node[lbl] at (0,-1.1) {пульт};
  \foreach \y/\n/\f in {1.5/{01 Quad-5in}/model01.yml,0/{02 Wing-900}/model02.yml,-1.5/{03 Cessna}/model03.yml}{
    \node[draw=etx,fill=white,rounded corners=2pt,minimum width=30mm,align=left,anchor=west] at (2.2,\y) {\n\\[-2pt]\scriptsize\ttfamily \f};
    \draw[arr,black!50] (1.4,0) -- (2.15,\y);}
  \begin{scope}[shift={(8.0,1.5)}]\pic[transform shape,scale=0.3] {quadtop};\end{scope}
  \begin{scope}[shift={(8.0,0)},scale=0.3]\path[airframe] (0,1) -- (-2.6,-0.9) -- (-2.4,-1.2) -- (0,-0.6) -- (2.4,-1.2) -- (2.6,-0.9) -- cycle;
     \path[surf] (-2.4,-1.2) -- (-1.1,-0.85) -- (-1.2,-0.6) -- (-2.5,-1.0) -- cycle;
     \path[surf] (2.4,-1.2) -- (1.1,-0.85) -- (1.2,-0.6) -- (2.5,-1.0) -- cycle;\end{scope}
  \begin{scope}[shift={(8.0,-1.5)}]\pic[transform shape,scale=0.3] {planetop};\end{scope}
  \foreach \y in {1.5,0,-1.5}{\draw[arr,black!50] (5.5,\y) -- (7.0,\y);}
\end{tikzpicture}
\caption{Каждая модель в пульте --- отдельный файл со всеми её настройками}
\end{figure}

\section{Список моделей}

Нажми \btn{MDL} --- откроется первая страница, \menu{Model Select}.

\begin{center}
\lcdpopup{MODEL SELECT}{1/13}{\lr{01 Quad-5in}{}\lsel{02 Wing-900}{}\lr{03 Cessna}{}\lr{04 ---}{}\lblank\lblank}%
{\linv{Select model}\\Create model\\Copy model\\Move model\\Delete model}
\end{center}

Долгое \btn{ENTER} на строке открывает меню:
\begin{itemize}
  \item \menu{Select model} --- сделать модель текущей (с ней будет работать пульт);
  \item \menu{Create model} --- создать новую;
  \item \menu{Copy model} --- скопировать (удобно: сделал одну модель, скопировал для второго
        такого же коптера);
  \item \menu{Move}, \menu{Delete} --- переместить, удалить.
\end{itemize}

При создании EdgeTX может предложить \textbf{мастер} (wizard): «самолёт», «крыло», «планер»,
«мультикоптер». Он задаст пару вопросов и сам настроит миксеры. Для первого раза это удобно,
но в книге мы разберём, как всё делается руками --- чтобы ты понимал, что происходит.

\section{Страница Setup}

\btn{PAGE>} от списка --- главная страница модели, \menu{Setup}.

\begin{center}
\lcdscreen{SETUP}{2/13}{\lr{Model name}{Quad-5in}\lr{Timer1}{ON SA↓ 4:00}\lr{ Minute call}{☑}\lr{ Countdown}{Voice}\lr{Timer2}{OFF}\lr{Extended trims}{☐}}
\end{center}

\subsection{Имя модели}

\menu{Model name} --- до 15 букв. Колесом выбираешь символ, \btn{ENTER} --- следующий,
долгое \btn{ENTER} --- заглавная/строчная. Давай понятные имена: \code{Quad-5in},
\code{Wing-900}, \code{Cessna}. Хорошая привычка --- наклеить такое же имя на саму модель.

\subsection{Таймеры}

У модели три таймера. Главный вопрос --- \emph{когда} таймер считает:

\begin{figure}[H]
\centering
\begin{tikzpicture}
\begin{axis}[width=13cm,height=6.2cm,xmin=0,xmax=10,ymin=-0.2,ymax=4.6,
  axis lines=left,xlabel={время},ylabel={},ytick=\empty,xtick=\empty,
  clip=false,every axis x label/.style={at={(axis description cs:1,0)},anchor=north east,font=\sffamily\scriptsize}]
  % газ
  \addplot[thick,black!70] coordinates {(0,3.6) (1,3.6) (1.5,4.4) (3,4.2) (4,3.6) (5.5,3.6) (6,4.5) (8,4.3) (8.5,3.6) (10,3.6)};
  \node[font=\sffamily\scriptsize,anchor=east] at (axis cs:0,4.0) {газ};
  \draw[black!30,dashed] (axis cs:0,3.6) -- (axis cs:10,3.6);
  % ON
  \addplot[line width=3pt,etx!60] coordinates {(0,2.6) (10,2.6)};
  \node[font=\sffamily\scriptsize,anchor=east] at (axis cs:0,2.6) {ON};
  % THs
  \addplot[line width=3pt,okgreen!70] coordinates {(1,1.8) (4,1.8)};
  \addplot[line width=3pt,okgreen!70] coordinates {(5.5,1.8) (8.5,1.8)};
  \node[font=\sffamily\scriptsize,anchor=east] at (axis cs:0,1.8) {THs};
  % TH%
  \addplot[line width=3pt,accent!80,opacity=0.35] coordinates {(1,1.0) (4,1.0)};
  \addplot[line width=3pt,accent!80,opacity=0.35] coordinates {(5.5,1.0) (8.5,1.0)};
  \addplot[line width=3pt,accent!80] coordinates {(1.5,1.0) (3,1.0)};
  \addplot[line width=3pt,accent!80] coordinates {(6,1.0) (8,1.0)};
  \node[font=\sffamily\scriptsize,anchor=east] at (axis cs:0,1.0) {TH\%};
  % THt
  \addplot[line width=3pt,warn!70] coordinates {(1,0.2) (10,0.2)};
  \node[font=\sffamily\scriptsize,anchor=east] at (axis cs:0,0.2) {THt};
\end{axis}
\end{tikzpicture}
\caption{Когда идёт счёт в разных режимах таймера (цветная полоса = таймер считает).
Сверху --- как двигался стик газа}
\end{figure}

\begin{center}\small
\renewcommand{\arraystretch}{1.25}
\begin{tabularx}{\linewidth}{@{}L{28mm}Y@{}}
\toprule
\textbf{Режим} & \textbf{Когда считает} \\
\midrule
\menu{OFF} & никогда (таймер выключен) \\
\menu{ON} & всегда, пока включён его переключатель \\
\menu{Start} & запускается переключателем и потом не останавливается \\
\menu{Throttle} (THs) & только пока газ не на минимуме \\
\menu{Throttle \%} (TH\%) & тем быстрее, чем больше газ (на полном газу --- как обычные часы) \\
\menu{Throttle Start} (THt) & запускается при первом движении газа и больше не останавливается \\
\bottomrule
\end{tabularx}
\end{center}

Ещё у таймера есть:
\begin{itemize}
  \item \menu{Switch} --- переключатель, который разрешает счёт (например, арм \sw{SA↓});
  \item \menu{Start} --- с какого времени начать: 0:00 --- считает вверх, 4:00 --- обратный отсчёт;
  \item \menu{Minute call} --- каждую минуту говорить голосом;
  \item \menu{Countdown} --- писк или голос в последние секунды;
  \item \menu{Persistent} --- не сбрасывать при выключении (для подсчёта общего налёта).
\end{itemize}

\begin{recipe}{таймер для квадрокоптера}
Timer1: \menu{ON}, \menu{Switch} = \sw{SA↓} (арм), \menu{Start} = \code{4:00},
\menu{Minute call} включён, \menu{Countdown} = \menu{Voice}.
Когда ты армишь коптер, пойдёт обратный отсчёт от 4 минут, пульт каждую минуту скажет, сколько
осталось, а в конце скажет, что время вышло. Пора садиться!
\end{recipe}

\subsection{Триммеры}

\begin{itemize}
  \item \menu{Extended trims} --- триммеры с увеличенным ходом (не нужен обычно).
  \item \menu{Trim step} --- насколько сдвигается триммер за одно нажатие: \menu{Fine}
        (мелко), \menu{Medium}, \menu{Coarse} (крупно), \menu{Exponential} (у нуля мелко,
        дальше крупнее).
  \item \menu{Display trims} --- показывать числа триммеров на главном экране.
\end{itemize}

\subsection{Газ}

\begin{itemize}
  \item \menu{Throttle source} --- что считать газом (обычно стик \sw{Thr}). От этого
        зависят таймеры и проверка газа при включении.
  \item \menu{Throttle trim idle only} --- триммер газа меняет только малые обороты
        (для моделей с бензиновым или калильным мотором).
\end{itemize}

\subsection{Предполётные проверки (Preflight checks)}

Это те самые проверки при включении, которые мы видели в главе~\ref{ch:radio}:

\begin{center}
\lcdscreen{SETUP}{2/13}{\lr{Preflight checks}{}\lr{ Throttle state}{☑}\lr{ Custom position}{☐}\lsel{ Switch positions}{SA↑ SB↑ SC↑}\lr{ Pots/Sliders}{OFF}\lblank}
\end{center}

\begin{enumerate}
  \item \menu{Throttle state} --- проверять, что газ внизу. Оставь включённым.
  \item \menu{Switch positions} --- поставь все переключатели так, как они должны стоять перед
        стартом (арм выключен!), встань на эту строку и нажми долгое \btn{ENTER} ---
        пульт запомнит положения.
\end{enumerate}

\begin{tip}
Обязательно добавь в проверку переключатель арма. Тогда пульт не включится с заармленным
коптером и не запустит моторы случайно.
\end{tip}

\subsection{Остальные пункты}

\begin{itemize}
  \item \menu{Internal RF} и \menu{External RF} --- радиомодули, следующая глава.
  \item \menu{Trainer} --- режим «учитель--ученик», глава~\ref{ch:trainer}.
  \item \menu{Use global functions} --- применять ли общие функции пульта к этой модели.
  \item \menu{Center beep} --- писк, когда стик или крутилка проходит через середину.
\end{itemize}

\begin{tryit}
\begin{enumerate}
  \item Создай новую модель, назови её \code{Test}.
  \item Настрой Timer1 на обратный отсчёт 3:00 по переключателю \sw{SA}.
  \item Включи минутные объявления и проверь, как пульт говорит время.
  \item Настрой предполётную проверку переключателей и попробуй включить пульт с \sw{SA} внизу.
\end{enumerate}
\end{tryit}

\begin{summary}
\begin{itemize}
  \item Одна настоящая модель --- одна модель в пульте.
  \item \btn{MDL} → список моделей; долгое \btn{ENTER} --- создать, копировать, удалить.
  \item На странице \menu{Setup}: имя, таймеры, триммеры, газ, проверки, радиомодуль.
  \item Предполётные проверки защищают от случайного запуска мотора.
\end{itemize}
\end{summary}
